%%%%%%%%%%%%%%%%%%%%%%%%%%%%%%%%%%%%%%%%%%%%%%%%%%%%%%%%%%%%
%%  This Beamer template was created by Cameron Bracken.
%%  Anyone can freely use or modify it for any purpose
%%  without attribution.
%%
%%  Last Modified: January 9, 2009
%%

\documentclass[xcolor=x11names,compress]{beamer}

%% General document %%%%%%%%%%%%%%%%%%%%%%%%%%%%%%%%%%
\usepackage{graphicx}
\usepackage{comment} % Allows commenting large blocks using \begin{comment} \end{comment} commands.
%\usepackage{enumitem}

%%% For Tikz Plots
\usepackage{xcolor} % For mixing colors.
\usepackage{tikz}
\usepackage{pgfplots}

%\pgfplotsset{          
%	axis line style={gray},
%         width=1\linewidth,    
%	height=0.35\linewidth 
%            }
% \usepackage{pgffor}

\usepackage{subcaption} 
\usepgfplotslibrary{fillbetween}

\usetikzlibrary{patterns}
\usetikzlibrary{arrows}
\usetikzlibrary{shapes,snakes} % For non-circle nodes

%%% Change Math font for product
\DeclareSymbolFont{replacement-operators}{OMX}{txex}{m}{n}

\DeclareMathSymbol{\bigprod}{\mathop}{replacement-operators}{"59}
\DeclareMathSymbol{\smallprod}{\mathop}{replacement-operators}{"51}
\renewcommand\prod{\mathop{\mathchoice
  {\bigprod}%
  {\smallprod}%
  {\smallprod}%
  {\smallprod}}}

%%%%%%%%%%%%%%%%%%%%%%%%%%%%%%%%%%%%%%%%%%%%%%%%%%%%%%

\usepackage[absolute,overlay]{textpos}
  \setlength{\TPHorizModule}{1mm}
  \setlength{\TPVertModule}{1mm}

\setbeamertemplate{navigation symbols}{} 
\setbeamertemplate{itemize items}[circle]
%% Beamer Layout %%%%%%%%%%%%%%%%%%%%%%%%%%%%%%%%%%
%\useoutertheme[subsection=false,shadow]{miniframes}
\useoutertheme[subsection=false,shadow]{}
\useinnertheme{default}
\usefonttheme{serif}
\usepackage{palatino, comment, color, mdsymbol, eucal, bbding}

% Game environments
\usepackage{sgamevar}
\renewcommand{\gamestretch}{2}

\setbeamerfont{title like}{shape=\scshape}
\setbeamerfont{frametitle}{shape=\scshape}

\setbeamercolor*{lower separation line head}{bg=DeepSkyBlue4} 
\setbeamercolor*{normal text}{fg=black,bg=white} 
\setbeamercolor*{alerted text}{fg=red} 
\setbeamercolor*{example text}{fg=black} 
\setbeamercolor*{structure}{fg=black} 
\setbeamercolor{frametitle}{fg=darkblue}
 
\setbeamercolor*{palette tertiary}{fg=black,bg=black!10} 
\setbeamercolor*{palette quaternary}{fg=black,bg=black!10} 



\renewcommand{\(}{\begin{columns}}
\renewcommand{\)}{\end{columns}}
\newcommand{\<}[1]{\begin{column}{#1}}
\renewcommand{\>}{\end{column}}

\let\inf\relax \DeclareMathOperator*\inf{\vphantom{p}inf}
\let\max\relax \DeclareMathOperator*\max{\vphantom{p}max}
\let\min\relax \DeclareMathOperator*\min{\vphantom{p}min}
\definecolor{darkblue}{rgb}{0, 0.08, 0.55}

%% Restoring caligraphic font to default
\DeclareMathAlphabet{\mathcal}{OMS}{cmsy}{m}{n}

%%Uniform distribution operator
%\DeclareMathOperator*{\argmax}{}
\newcommand{\argmax}{\mathrm{argmax}}
\newcommand{\argmin}{\mathrm{argmin}}


%% General Macros.
\newcommand{\samplepoints}{samples=50}
\newcommand{\meddef}{samples=200}
\newcommand{\highdef}{samples=500}


\newcommand{\co}{\text{co}}
\newcommand{\ca}{\text{ca}}
\newcommand{\ext}{\text{ext}}
\newcommand{\marg}{\text{marg}}
\newcommand{\ddd}{\text{ d}}
\newcommand{\dd}{\text{d}}
\newcommand{\supp}{\text{supp}}
\newcommand{\bbE}{\mathbb E}
\newcommand{\dom}{\text{dom}}
\newcommand{\inter}{\text{int}}
\newcommand{\cl}{\text{cl}}

\newcommand{\eff}{\alpha}
\newcommand{\effb}{\beta}
\newcommand{\effc}{\gamma}
\newcommand{\Eff}{\mathcal{A}}
\newcommand{\minx}{L}%{\underline{x}}
\newcommand{\maxx}{H}%{\bar{x}}
\newcommand{\waged}{\chi}
\newcommand{\uinv}{n}

\newcommand{\acomment}[1]{}

%%Divide slide into four command
\newcommand\FourQuad[4]{%
    \begin{minipage}[b][.35\textheight][t]{.47\textwidth}#1\end{minipage}\hfill%
    \begin{minipage}[b][.35\textheight][t]{.47\textwidth}#2\end{minipage}\\[0.5em]
    \begin{minipage}[b][.35\textheight][t]{.47\textwidth}#3\end{minipage}\hfill
    \begin{minipage}[b][.35\textheight][t]{.47\textwidth}#4\end{minipage}%
}
\newcommand\HalfandTwoQuad[3]{%
    \begin{minipage}[b][.35\textheight][t]{.47\textwidth}#1\end{minipage}\hfill%
    \begin{minipage}[b][.35\textheight][t]{.47\textwidth}#2\end{minipage}\\[0.5em]
    \begin{minipage}[b][.35\textheight][t]{1\textwidth}#3\end{minipage}%
}



%%%%%%%%%%%%%%%%%%%%%%%%%%%%%%%%%%%%%%%%%%%%%%%%%%




\begin{document}
\begin{frame}
\vspace{-0.75cm}
	\title{{\color{darkblue}\textit{Flexible Moral Hazard Problems}}}
	\subtitle{}

	\vspace{1cm}

\author{
%\\ \bigskip
\begin{minipage}{0.4\textwidth}\centering  
George Georgiadis\\ \smallskip \centering \footnotesize Northwestern University
\end{minipage} 
\begin{minipage}{0.4\textwidth}\centering  
Doron Ravid\\ \smallskip \centering \footnotesize University of Chicago 
\end{minipage}
\\ \bigskip
\begin{minipage}{0.4\textwidth}\centering  
Bal\'azs Szentes \\ \smallskip \centering \footnotesize University of Hong Kong
\end{minipage}
}



\bigskip
	\date{
		December 2023
	}
\titlepage
\end{frame}


\begin{frame}{The Informativeness Principal}
Classic moral hazard: effort is binary or an interval. \\ \medskip
\textbf{Main result:} Contracts are motivated by informativeness. \\ \medskip
\onslide<2->{
This \textcolor{blue}{\emph{informativeness principle}} is a key tenant of contract theory:\\ \medskip
}

\onslide<2->{
\begin{quote}
``In the late 1970s, Bengt Holmström demonstrated how a principal (e.g., a company’s shareholders) should design an optimal contract for an agent (the company’s CEO), whose action is partly unobserved by the principal. Holmström’s \textcolor{blue}{informativeness principle} stated precisely how this contract should link the agent’s pay to performance-relevant information. Using the basic principal-agent model, he showed how the optimal contract carefully weighs risks against incentives. In later work, Holmström generalised these results to more realistic settings...''
\end{quote}
\footnotesize{(The Swedish Academy's 2016 nobel prize press release)}
}

\end{frame}


\begin{frame}{Informativeness Principle: The Issues}

Two implications of informativeness principle:\medskip
\begin{enumerate}
\item Need strong assumptions for wage to increase in output.\medskip
\item Noisier measures $\implies$ lower powered incentives. \bigskip
\end{enumerate}

\onslide<2->{
These implications are empirically problematic:\medskip
\begin{itemize}
\item Contracts almost always increasing in output. \medskip
\item Incentives weakly connected to noise (Prendergast 1999). \bigskip
\end{itemize}
}

\onslide<3->{Innes (1990): Get monotonicity from free money burning.}

\end{frame}

\begin{frame}{This Paper: Flexible Output Choice}
We allow agent to choose \emph{any} output distribution. \\ \bigskip

Main results: \medskip
\begin{itemize}
\item Contracts motivated by agent's marginal costs.\medskip
\item FOSD monotone costs $\implies$ increasing contracts. \medskip
\item Result by omission: Informativeness plays no role.
\end{itemize}
\end{frame}

\begin{frame}
	\frametitle{}
	\begin{center}
		{\color{darkblue} {\huge Two Examples}}
	\end{center}
\end{frame}

\begin{frame}{Common Setup for Examples}

A principal (she) contracts with an agent (he).\medskip

\begin{itemize}

\item Compact set $X\subset \mathbb{R}$ of possible outputs.\medskip

\item Principal offers agent a (bounded) contract: $w:X \rightarrow \mathbb{R}$.\medskip

\item Agent can opt out and get $u_0$.\medskip

\item If opts in, agent covertly chooses $\eff \in \Eff \subseteq \Delta (X)$.\medskip

\item Effort costs: $C:\Eff \rightarrow \mathbb{R}_{+}$, increasing in FOSD. \medskip

\item Payoffs: 
\[
\text{Principal: } x - w \quad \quad \text{Agent: } u(w) - C(\eff).
\]
$u$: strictly increasing, differentiable, unbounded, concave.
\end{itemize}

\end{frame}

\begin{frame}{Standard Binary Effort Model}
\begin{textblock}{100}(15,13)
\[ 
X = \left[\minx,\maxx\right], \quad \quad \Eff = \{\eff_l,\eff_h\}.
\]
\end{textblock}
\begin{textblock}{100}(11,25)
\only<2->{
Suppose principal wants to implement $\eff_h$. \\ \bigskip
}
\only<3->{
Then she offers a contract $w$ that solves:
}
\only<3>{
\begin{align}
\min_{w(\cdot)} & \int w(x) \eff_h(\dd x)  \notag \\ 
\text{s.t.} 
& \int u\circ w(x) (\eff_h - \eff_l) (\dd x) \geq C(\eff_h) - C(\eff_l) \tag{IC}\\
& \int u\circ w(x) \eff_h (\dd x) - C(\eff_h)\geq u_0. \tag{IR}
\end{align}
}
\only<4->{
\[
\min_{w(\cdot)} \int w(x) \eff_h(\dd x) \quad \text{s.t.} \quad \text{(IC) and (IR)}.
\]
}
\only<5->{
The FOC from this cost minimization problem is:
\[
\frac{1}{u^{\prime }\left( w\left( x \right) \right) }=\lambda +\mu \left[
1-\frac{f_l\left( x \right) }{f_h\left( x \right) }\right]
\]
}
\only<6->{
So: $w$ is monotone $\iff$ MLRP holds.
}
\end{textblock}
\end{frame}

\begin{frame}{Flexible Binary Output Model}
\begin{textblock}{100}(15,13)
\[
X = \textcolor{blue}{ \left\{ \minx,\maxx \right\}}, \quad \quad \Eff = \textcolor{blue}{\Delta X \equiv [0,1].}
\]
\end{textblock}
\begin{textblock}{100}(11,25)
\only<2->{
Suppose also $C$ is convex and differentiable.\\ \medskip
}
\only<3-4>{
Agent's problem is:
\[
\max_{\eff \in [0,1]} [\eff u \circ w(\maxx) + (1-\eff) u \circ w(\minx) - C(\eff)].
\]
}
\only<4->{
Agent's FOC for choosing $\eff \in (0,1)$:
\[
u\circ w(\maxx) - u \circ w(\minx) = C'(\eff) \only<5->{\onslide<6->{\iff w(\maxx) = u^{-1}(u\circ w(\minx) + C'(\eff)).}}
\]
\onslide<7->{
Implications:\medskip}
\begin{itemize}
\item<8-> Cost minimization is trivial: min $w(\minx)$ s.t. IR.\medskip
\item<9-> Shape of contract determined by $C'$ and $u$. \medskip
\item<10-> IC contracts are monotone: 
\[
w(\maxx) = u^{-1}(u\circ w(\minx) + C'(\eff)) \geq u^{-1}(u\circ w(\minx)) = w(\minx).
\]
\end{itemize}
}
\end{textblock}
\end{frame}


\begin{frame}{Our Main Model}

A principal (she) contracts with an agent (he).\medskip

\begin{itemize}

\item Compact set $X\subset \mathbb{R}$ of possible outputs.\medskip

\item Principal offers agent a (bounded) contract: $w:X \rightarrow \mathbb{R}_{+}$.\medskip

\item \textcolor{blue}{Limited liability: $w(\cdot) \geq 0$.}\medskip

\item Agent covertly chooses $\textcolor{blue}{\eff \in \Eff = \Delta (X)}$.\medskip

\item Effort costs: $C:\Eff \rightarrow \mathbb{R}_{+}$, continuous, increasing in FOSD. \medskip

\item Payoffs: 
\[
\text{Principal: } x - w \quad \quad \text{Agent: } u(w) - C(\eff),
\]
\textcolor{blue}{
$u$: increasing, continuous, unbounded \& $u(0)=0$.
}
\end{itemize}

\end{frame}

\begin{frame}{Assumptions on the Cost}

\onslide<2->{
\noindent \textbf{Without loss:} $C$ is convex. \\ \medskip
(if not, replace $\eff$ with cheapest mixing that averages to $\eff$) \\ \bigskip
}
\onslide<3->{
\noindent \textbf{Assumption. (smoothness) }$C$ is \textbf{Gateaux
differentiable}: every $\eff$ admits a bounded $k_{\eff}:X\rightarrow \mathbb{R}$ s.t. 
\begin{equation*}
\lim_{\epsilon \downarrow 0}\frac{1}{\epsilon }\left[ C(\eff +\epsilon (\effb
 -\eff ))-C(\eff )\right] =\int k_{\eff}\left( x\right) (\effb-\eff)\left( \dd x\right)
\end{equation*}%
for all $\effb\in \Eff$. \\ \bigskip
}
\onslide<4->{
\textbf{Notes:}\medskip
\begin{itemize}
\item<4-> $k_\alpha(x)$: MC of increasing probability of output $x$. \medskip
\item<5-> Without loss of generality: Normalize $\inf k_{\eff}(X) = 0$. 
\end{itemize}
}
\end{frame}

\begin{frame}{Example 1: Finite Outputs}
Suppose $X=\{x_1,\ldots,x_N\}$. \\ \bigskip
Think of $\eff = (\eff_1,\ldots,\eff_N)$ as an element of $\mathbb{R}^{N}$.\\ \bigskip
Then Gateaux differentiability $\iff$ differentiable. \\ \bigskip
In this case,
\[
k_\eff (x_i) = \frac{\partial C}{\partial \alpha_i} (\eff)  \onslide<2->{- \min_{j=1,\ldots, N}\frac{\partial C}{\partial \alpha_j} (\eff).}
\]
\onslide<3->{
In this case: $C$ convex $\implies$ $C$ differentiable Lebesgue a.e.}
\end{frame}


\begin{frame}{Example 2: Affine Costs}
Suppose we have an increasing and bounded function,
\[
k: X \rightarrow \mathbb{R}_{+},
\]
with $k(\min \ X) = 0$, and costs are given by
\[
C(\eff) = \int k(x) \eff(\dd x).
\]
This cost is convex and differentiable, with derivative
\[
k_{\eff} (\cdot) = k(\cdot).
\]
\end{frame}

\begin{frame}{Example 3: Transformations of Affine Costs}
Consider $k$ as before, and let
\[
\Psi: \mathbb{R}\rightarrow \mathbb{R}
\]
be increasing, convex, differentiable, and $\Psi(0)=0$. Then,
\[
C(\eff) = \Psi \left( \int k(x) \eff (\dd x) \right)
\]
satisfies our assumptions. The derivative is
\[
k_{\eff}(\cdot) = \Psi'\left(\int k(x) \eff (\dd x) \right) k(\cdot).
\]
\end{frame}

\begin{frame}
	\frametitle{}
	\begin{center}
		{\color{darkblue} {\huge Cost Minimization}}
	\end{center}
\end{frame}

\begin{frame}{First-Order Approach}
\textbf{Lemma.} For a bounded $v:X\rightarrow \mathbb{R}$, and $\eff \in 
\Eff$, 
\begin{equation*}
\eff \in \arg \max_{\effb\in \Eff}\left[ \int v(x)\effb\left( \dd x\right) -\textcolor{blue}{C(\effb)}\right]
\end{equation*}%
if and only if 
\begin{equation*}
\eff \in \arg \max_{\effb \in \Eff}\left[ \int v(x) \effb(\dd x) - \textcolor{blue}{\int k_{\eff
}(x) \effb (\dd x)} \right]
\end{equation*}%
(the ``only if'' direction also works if $C$ is not convex)
\end{frame}

\begin{frame}{Relationship to Standard FOC}
Consider the problem:
\[
\max_{a \in [0,1]} \left[av - c(a)\right]
\]
where $v \in \mathbb{R}$ and $c$ is convex and differentiable. \\ \medskip
\onslide<2->{
Standard way of writing FOC for optimal $a^* \in (0,1)$ is
\[
v - c'(a^*)=0.
\]
}% 
\onslide<3->{%
An equivalent way of writing the above condition is:
\[
a^* \in \argmax_{a \in [0,1]} [av - ac'(a^*)].
\]
The lemma generalizes the second formulation. 
}
\end{frame}

\begin{frame}{First-Order Approach}
\textbf{Lemma.} For a bounded $v:X\rightarrow \mathbb{R}$, and $\eff \in 
\Eff$, 
\begin{equation*}
\eff \in \arg \max_{\effb\in \Eff}\left[ \int v(x)\effb\left( \dd x\right) -\textcolor{blue}{C(\effb)}\right]
\end{equation*}%
if and only if 
\begin{equation*}
\eff \in \arg \max_{\effb \in \Eff}\left[ \int v(x) \effb(\dd x) - \textcolor{blue}{\int k_{\eff
}(x) \effb (\dd x)} \right]
\end{equation*}%
(the ``only if'' direction also works if $C$ is not convex)
\end{frame}

\begin{frame}{Characterization of IC}
Say a contract-distribution pair $(w,\eff)$ is \textbf{IC} if
\[
\eff \in \argmax_{\effb \in \Eff} \left[\int u\circ w(x) \effb(\dd x) - C(\effb)\right].
\]

\onslide<2->{
\textbf{Proposition.} $(w,\eff)$ is IC if and only if a $m \in \mathbb{R}$ exists such that
\[
w(x) \leq u^{-1}(k_\eff(x) + m)
\]
for all $x$, and with equality $\eff$-almost surely.
}
\end{frame}


\begin{frame}{}
\begin{textblock}{105}(11,10)
\textbf{Proposition.} $(w,\eff)$ is IC if and only if a $m \in \mathbb{R}$ exists such that
\[
w(x) \leq u^{-1}(k_\eff(x) + \textcolor<5->{blue}{m})
\]
for all $x$, and with equality $\eff$-almost surely.
\end{textblock}
\begin{textblock}{105}(11,40)
\textbf{Proof.} \onslide<2->{
By Lemma, $(w,\eff)$ is IC if and only if $\eff$ solves
\[
\max_{\effb \in \Delta X}\int \left[ u\circ w (x) - k_\eff(x)\right] \effb(\dd x),
\]
}

\onslide<3->{
\noindent or equivalently, the following holds $\eff$-almost surely:
\[
\begin{split}
u\circ w (x) - k_\eff(x) & = \sup (u\circ w - k_\eff )(X) 
\onslide<4->{\\ \iff  w(x) & = u^{-1} \left( k_\eff(x) + \textcolor<5->{blue}{\sup (u\circ w - k_\eff )(X)} \right).}
\end{split}
\]
}
\end{textblock}
\end{frame}

\begin{frame}{}
\begin{center}
\end{center}
\begin{textblock}{105}(11,10)
\textbf{Proposition.} $(w,\eff)$ is IC if and only if a $m \in \mathbb{R}$ exists such that
\[
w(x) \leq \textcolor{blue}{u^{-1}(k_\eff(x) + m) =: w_{m,\alpha}(x)}
\]
for all $x$, and with equality $\eff$-almost surely.
\end{textblock}

\begin{textblock}{105}(11,35)
\only<2-4>{
\begin{figure}
\centering{}    \begin{tikzpicture}[scale=0.8]
			\begin{axis} 
				[
					\samplepoints,
%					ticks=none,
					ymax=1.1,
					ymin=-0.03,
					xmin=0,
					xmax=1.2,
					axis on top=false,
					axis x line = middle,
%					axis x line = none,
%					axis line style={-},
%					axis y line = none,
%					every inner x axis line/.append style={|-|},
					axis y line = middle,
					axis line style={gray},	
					xlabel=$x$,
					ylabel=$ $,
%					every axis x label/.style= {at={(ticklabel cs:0.5)},anchor=center},					
%				    xtick={0,1},
				    xtick= \empty,				    	
%				    xticklabels={$0$,$1$},
%				    ytick={0,1},
%				    yticklabels={$0$,$1$}
				    ytick=\empty,
%				    yticklabels={}
				]
				
\only<1-3>{
	\addplot [blue, line width=2, domain=0:1] {2*(1/8 + (x - 0.5)^3)} node[right]{$w_{m_1,\eff}$};

	\addplot [blue, line width=2, domain=0:1] {2*(2/8 + (x - 0.5)^3)} node[right]{$w_{m_2,\eff}$};
	

	\addplot [blue, line width=2, domain=0:1] {2*(3/8 + (x - 0.5)^3)} node[right]{$w_{m_3,\eff}$};
}

\only<4->{
	\addplot [blue, line width=2, domain=0:1] {2*(2/8 + (x - 0.5)^3)} node[right]{$w_{m,\eff}$};
	}
	
	
%	\addplot+ [no marks, red, line width=2] coordinates {(0,0) (1/4,15/32) (5/8,15/32) (3/4,17/32) (1,17/32)};

\only<3->{
\addplot+ [no marks, solid, violet, line width=2] coordinates {(0.25,0) (0.75,0)} node[above]{$\supp (\eff)$};
}

\only<3->{
	
	\addplot+ [no marks, red, solid, line width=2] coordinates {(3/4,17/32) (1,20/32)} node[right] {$w$};

	\addplot [red, line width=2, dashed, domain=0.25:0.75] {2*(2/8 + (x - 0.5)^3)};
	
	\addplot+ [no marks, red, solid, line width=2] coordinates {(0,0) (1/4,15/32)};
}
	\end{axis}
\end{tikzpicture}
\end{figure} 

}

\only<5->{
\begin{columns}
\begin{column}{0.5\textwidth}
\begin{figure}
\centering{}    \begin{tikzpicture}[scale=0.8]
			\begin{axis} 
				[
					\samplepoints,
%					ticks=none,
					ymax=1.1,
					ymin=-0.03,
					xmin=0,
					xmax=1.2,
					axis on top=false,
					axis x line = middle,
%					axis x line = none,
%					axis line style={-},
%					axis y line = none,
%					every inner x axis line/.append style={|-|},
					axis y line = middle,
					axis line style={gray},	
					xlabel=$x$,
					ylabel=$ $,
%					every axis x label/.style= {at={(ticklabel cs:0.5)},anchor=center},					
%				    xtick={0,1},
				    xtick= \empty,				    	
%				    xticklabels={$0$,$1$},
%				    ytick={0,1},
%				    yticklabels={$0$,$1$}
				    ytick=\empty,
%				    yticklabels={}
				]
				

\only<10->{
	\addplot [blue, line width=2, domain=0:1] {2*(1/8 + (x - 0.5)^3)} node[right]{$\textcolor{violet}{w_{m^{*}_{\eff},\eff}}$};
	}
	

\only<1-6>{
	\addplot [blue, line width=2, domain=0:1] {2*(2/8 + (x - 0.5)^3)} node[right]{$w_{m,\eff}$};
	}
\only<7-9>{
	\addplot [blue, line width=2, domain=0:1] {2*(2/8 + (x - 0.5)^3)} node[above right]{$w_{m,\eff}$}
};


%	\addplot+ [no marks, red, line width=2] coordinates {(0,0) (1/4,15/32) (5/8,15/32) (3/4,17/32) (1,17/32)};

\addplot+ [no marks, solid, violet, line width=2] coordinates {(0.25,0) (0.75,0)} node[above]{$\supp (\eff)$};

\only<4-6>{
	
	\addplot+ [no marks, red, solid, line width=2] coordinates {(3/4,17/32) (1,20/32)} node[right] {$w$};

	\addplot [red, line width=2, dashed, domain=0.25:0.75] {2*(2/8 + (x - 0.5)^3)};
	
	\addplot+ [no marks, red, solid, line width=2] coordinates {(0,0) (1/4,15/32)};
}
	
\only<7-9>{
	\addplot [red, line width=2, dashed, domain=0:1] {2*(2/8 + (x - 0.5)^3)} node[below right]{$=w$};
}

\only<10->{
	\addplot [red, line width=2, dashed, domain=0:1] {2*(1/8 + (x - 0.5)^3)};% node[below right]{$=w$};
}
	
	
%	\node [circle, fill=violet, scale=0.7] at (axis cs:3/8,0){};


			\end{axis}
\end{tikzpicture}
\end{figure} 

\end{column}
\begin{column}{0.5\textwidth}
\onslide<5->{
\textbf{Implications:} \medskip
\begin{enumerate}[(i)]
\item<6-> wlog: set $w=w_{m,\eff}$, optimize $m$. 
\item<8-> Every $\eff$ is implementable.
\item<9-> Cost minimizing $m$ is:
\[
m^{*}_{\eff} = - \inf k_{\eff}(X) = 0.
\]

\end{enumerate}
}
\end{column}
\end{columns}
}

\end{textblock}

\end{frame}


\begin{frame}{Monotone Contracts}
\textbf{Claim.} $C$ is FOSD monotone $\iff$ $w_{m,\eff}$ is increasing $\forall \eff, m$. \\ \bigskip

\onslide<2->{
(\textbf{Fact.} $C$ is FOSD monotone $\iff$ $k_{\eff}$ is increasing $\forall \eff$).\\ \bigskip
}

\onslide<3->{
\textbf{Proof of Claim.} 
Recall,
\[
w_{m,\eff}(x) = u^{-1}( k_{\alpha}(x) + m).
\]
}

\onslide<4->{
Therefore,
\[
\begin{split}
w_{m,\eff} \text{ is increasing }\forall \eff,m  
& \iff k_{\eff} \text{ is increasing } \forall \eff
\onslide<5->{\\ & \iff C \text{ is FOSD monotone}.}
\end{split}
\]
}

\onslide<4->{\textbf{Explanation: }}\only<4>{because $u^{-1}$ is increasing.}\only<5>{by the Fact .}

\end{frame}




\begin{frame}{}
\textbf{Proposition.} $(w,\eff)$ is IC if and only if a $m \in \mathbb{R}$ exists such that
\[
w(x) \leq u^{-1}(k_\eff(x) + m) =: w_{m,\eff}(x)
\]
for all $x$, and with equality $\eff$-almost surely. \\ \bigskip


\textbf{Implications:} \medskip
\begin{enumerate}[(i)]
\item Without loss for principal to offer $w_{m,\eff}$ for some $m$. \medskip
\item A cheapest contract implementing $\eff$ is $w_{m^{*}_{\eff},\eff}$ for 
\[
m^{*}_{\eff} = - \inf k_{\eff}(X) = 0.
\]
\item FOSD monotone $C \implies$  wage is increasing without loss.
\end{enumerate}


\end{frame}

\begin{frame}{Example 2: Affine Costs}
Suppose we have an increasing and bounded function,
\[
k: X \rightarrow \mathbb{R}_{+},
\]
with $k(\inf \ X) = 0$, and costs are given by
\[
C(\eff) = \int k(x) \eff(\dd x).
\]
This cost is convex and differentiable, with derivative
\[
k_{\eff} (\cdot) = k(\cdot).
\]
\end{frame}

\begin{frame}{Example 2: Affine Costs (continued)}
In this example, $(w,\eff)$ is IC if and only if an $m$ exists s.t.
\[
w(x) \leq w_{m,\eff}(x)= u^{-1}(k(x) + m),
\]
for all $x$, with equality $\eff$-almost surely.\\ \bigskip
\onslide<2->{
Thus, cost minimizing wage for \emph{any} $\eff$ is
\[
w^{*}(\cdot) = u^{-1}\circ k(\cdot).
\]
}
\end{frame}

\begin{frame}{Example 2: Affine Costs (continued)}
\begin{itemize}
\item<1-> Notice agent's utility from any $\eff$ is zero:
\[
\begin{split}
\onslide<2->{\int u\circ w^{*}(x) \eff (\dd x) - C(\eff)}
\onslide<3->{& = \int u\circ u^{-1}(k(x)) \eff (\dd x) - \int k(x) \eff (\dd x) }
\onslide<4->{\\ & = \int k(x) \eff (\dd x) - \int k(x) \eff (\dd x) =0.}
 \end{split}
\]

\item<5-> In this case, a deterministic output are optimal. \\ \medskip
\onslide<6->{\textbf{Explanation:} Principal's problem is:
\[
\max_{\eff \in \Delta X} \int \left[x - w^*(x)\right]\eff(\dd x).
\]
}
\onslide<7->{So, optimal to implement a deterministic output $x^*$ s.t.
\[
x^* \in \argmax_{x \in X}\left[x - w^*(x))\right].
\]
}
\end{itemize}
\end{frame}

\begin{frame}{Example 3: Transformations of Affine}
Consider $k$ as before, and let
\[
\Psi: \mathbb{R}_{+}\rightarrow \mathbb{R}_{+}
\]
be increasing, convex, differentiable, and $\Psi(0)=0$. Then,
\[
C(\eff) = \Psi \left( \int k(x) \eff (\dd x) \right)
\]
satisfies our assumptions. The derivative is
\[
k_{\eff}(\cdot) = \Psi'\left(\int k(x) \eff (\dd x) \right) k(\cdot).
\]
\end{frame}

\begin{frame}{Example 3: Transformations of Affine (cont.)}
In this example, $(w,\eff)$ is IC if and only if an $m$ exists s.t.
\[
w(x) \leq w_{m,\eff}(x)= u^{-1}\left(\Psi'\left(\int k(y) \eff (\dd y) \right) k(x) + m\right),
\]
for all $x$, with equality $\eff$-almost surely.\\ \bigskip
\onslide<2->{
Given $\eff$, a cost minimizing wage is
\[
w_{0,\eff}(\cdot) = u^{-1}\left(\Psi'\left(\int k(y) \eff (\dd y) \right) k(\cdot)\right).
\]
Wage depends on $\eff$, so principal's problem more complex. 
}
\end{frame}

\begin{frame}
	\frametitle{}
	\begin{center}
		{\color{darkblue} {\huge Profit Maximization}}
	\end{center}
\end{frame}

\begin{frame}{The Principal's Problem}
Let $w_{\eff}:=w_{0,\eff}$ be the cost minimizing wage implementing $\eff$. \\ \bigskip

The principal's problem is:
\[
\max_{\eff \in \Delta X} \left[\int  x \eff (\dd x)  - \int w_{\eff}(x) \eff (\dd x)\right].
\]
\end{frame}

\begin{frame}{Additional Assumptions}
\textbf{Nice Agent's Payoffs.} $\uinv:=u^{-1}$ is continuously differentiable.\\ \bigskip

\onslide<2->{
\textbf{Continuous Derivative.} $k_{\eff}$ is continuous, and $\eff \mapsto k_{\eff}$ is weak*-supnorm continuous. \\ \bigskip
}

\onslide<3->{
\textbf{2nd Order Differentiability.} Every $\eff$ admits a continuous function $h_{\eff}: X \times X \rightarrow \mathbb{R}$ such that for all $\effb$,
\[
\lim_{\epsilon\rightarrow 0} \frac{1}{\epsilon} \left[k_{\eff+\epsilon(\effb - \eff)}(\cdot) - k_{\eff}(\cdot) \right] = \int h_{\eff}(x,\cdot) (\effb - \eff)(\dd x),
\]
where convergence is in the supnorm.\\ \bigskip
}
\onslide<4->{For finite $X$: last two assumptions$\iff$twice differentiability.}
\end{frame}

\begin{frame}{Principal First Order Condition}
Define the function $\waged:X\rightarrow\mathbb{R}$,
\[
\waged_{\eff} (x) = \int h_{\eff}(x,y) (\uinv ' \circ k_{\eff})(y) \eff(\dd y).
\]

\onslide<2->{
\textbf{Theorem.} \\ \medskip A profit maximizing $\eff^*$ exists. Moreover, $\eff^*$ must solve
\[
\max_{\eff \in \Delta X} \int \left[x - w_{\eff^*}(x) - \waged_{\eff^*}(x)\right] \eff(\dd x).
\]
}
\end{frame}

\begin{frame}{Nice Profit Maximizing Distributions}
For every $\eff$, let 
\[
\pi_{\eff}(x) := x - w_{\eff}(x) - \waged_{\eff}(x).
\]

\onslide<2->{
\textbf{Corollary.} Suppose $X = [\minx,\maxx]$ and $\eff^*$ maximizes profits. Then,\medskip
\begin{enumerate}[(i)]
\item<3-> If $\pi_{\eff}$ is strictly \emph{quasiconcave} $\forall \eff$, then $|\supp \ \eff^*|=1$. \medskip

\item<4-> If $\pi_{\eff}$ is strictly \emph{quasiconvex} $\forall \eff$, then $\supp \ \eff^*\subseteq\{\minx,\maxx\}.$ \medskip

\item<5-> If $w_\eff + \waged_\eff$ is non-affine analytic $\forall\eff$, $\supp(\eff^*)$ is countable. \bigskip

\end{enumerate}
}

\end{frame}

\begin{frame}{Discrete Distributions and Bonuses}

\begin{textblock}{105}(11,15)
\textbf{Remark.} Suppose $X=[\minx,\maxx]$, $C$ is FOSD monotone, and \\ \medskip $\supp(\eff)$ is countable. Then,
\[
w(x) = w_{\eff}\left(\max\left\{y \in \supp (\eff) \cup \{\minx\}: y \leq x\right\}\right)
\]
is a cost minimizing contract that implements $\eff$.
\end{textblock}

\begin{textblock}{105}(11,45)
\begin{figure}
\centering{}    \begin{tikzpicture}[scale=0.8]
			\begin{axis} 
				[
					\samplepoints,
%					ticks=none,
					ymax=0.6,
					ymin=-0.03,
					xmin=-0.03,
					xmax=1.1,
					axis on top=false,
					axis x line = middle,
%					axis x line = none,
%					axis line style={-},
%					axis y line = none,
%					every inner x axis line/.append style={|-|},
					axis y line = middle,
					axis line style={gray},	
					xlabel=$x$,
					ylabel=$ $,
%					every axis x label/.style= {at={(ticklabel cs:0.5)},anchor=center},					
%				    xtick={0,1},
				    xtick= \empty,				    	
%				    xticklabels={$0$,$1$},
%				    ytick={0,1},
%				    yticklabels={$0$,$1$}
				    ytick=\empty,
%				    yticklabels={}
				]
				

	\addplot [blue, line width=2, domain=0:1] {2*(1/8 + (x - 0.5)^3)} node[right]{$w_{\eff}$};


	\addplot+ [no marks, red, solid, line width=2] coordinates {(0,0) (1/4,0) (1/4,7/32) (3/4,7/32) (3/4,9/32) (1,9/32)} node[right]{$w$};
	
	\node [circle, fill=violet, scale=0.7] at (axis cs:0,0){};
	\node [circle, fill=violet, scale=0.7] at (axis cs:0.25,0){};
	\node [circle, fill=violet, scale=0.7] at (axis cs:0.75,0){};
	\node[no marks, violet] at (axis cs:0.6,0.03){$\supp (\eff)$};
	\addplot+ [no marks, violet, dashed, line width=2] coordinates {(3/4,0) (3/4,9/32)};

	
	\end{axis}
\end{tikzpicture}
\end{figure} 


\end{textblock}

\end{frame}


\begin{frame}{Example 3: Transformations of Affine (cont.)}
Consider again the cost function
\[
C(\eff) = \Psi \left( \int k(x) \eff (\dd x) \right),
\]
and suppose $\Psi$ is twice continuously differentiable. \\ \bigskip
Then,
\[
h_{\eff}(x,y) = \Psi''\left( \int k(z) \eff (\dd z)\right) k(x)k(y).
\]
\onslide<2->{
So,
\[
\waged_{\eff}(x) = \textcolor{gray}{\left[\int \Psi''\left( \int k(z) \eff (\dd z)\right) \uinv'\left(\Psi'\left( \int k(z) \eff (\dd z)\right) \right)(k(y))^2 \eff(\dd y)\right]}k(x).
\]
}
\end{frame}

\begin{frame}{Example 3: Transformations of Affine (cont.)}
Hence, the Theorem implies $\eff^*$ is optimal only if 
\[
\supp \ \eff^* \subseteq \argmax_{x \in X} \left[x - w_{\eff^*}(x) - \waged_{\eff^*}(x)\right],
\]
where \bigskip
\begin{itemize}
\item $w_{\eff^*}(x) = \uinv\bigg(\Psi'\left( \int k(z) \eff^*(\dd z) \right)k(x) \bigg)$, \bigskip
\item $\waged_{\eff^*}(x) = \left[\text{positive stuff that depends on }\eff^*\right] k(x). $
\end{itemize}
\end{frame}

\begin{frame}{Specializing Example 3}
Suppose:
\[
X=[0,1],\quad 
u(y) = y^\rho, \quad 
k(x) = x^{\gamma}\quad 
\Psi(a) = \frac{a^{1+\lambda}}{1+\lambda},
\]
where $\gamma, \lambda >0$, and $\rho \in (0,1)$. \onslide<2->{Then,
\[
w_{\eff}(x) = [\text{positive stuff w/o }x ] \ x^{\frac{\gamma}{\rho}}, \quad 
\waged_{\eff}(x) = [\text{positive stuff w/o }x ] \ x^{\gamma}.
\]
}

\onslide<3->{
So for non-degenerate $\eff$, $\pi_{\eff}(x)=x - w_{\eff}(x) - \waged_{\eff}(x)$ is \medskip
\begin{itemize}
\item<4-> strictly concave if $\gamma \geq 1 (>\rho)$.\medskip
\item<5-> strictly convex if $\gamma \leq \rho$. \medskip
\item<6-> non-constant and analytic for all parameters.
\end{itemize}
}

\end{frame}

\begin{frame}{Specializing Example 3: Concave Case}
Suppose $\gamma \geq 1$, so $\pi_{\eff}$ is strictly concave. \\ \bigskip
By Corollary: only deterministic distributions can be optimal. \\ \bigskip
\onslide<1->{
For $x \in X=[0,1]$, let $\delta_{x}$ be deterministic distribution yielding $x$.\\ \bigskip
}

\onslide<1->{
Calculating $w_{\delta_{x}}(x) = x^{(1+\lambda)\frac{\gamma}{\rho}}$, we get $x^*$ is optimal iff
\[
\max_{x\in X} \left[x - x^{(1+\lambda)\frac{\gamma}{\rho}}\right]
\]
}

\onslide<1->{
Objective is concave, so can use FOC to get solution:
\[
x^* = \left[(1+\lambda)\frac{\gamma}{\rho}\right]^{\frac{1}{1-(1+\lambda)\frac{\gamma}{\rho}}}.
\]
}
\end{frame}


\begin{frame}{Specializing Example 3: Convex Case}
Suppose $\gamma \leq \rho$, so $\pi_{\eff}$ is strictly convex.\\ \bigskip
By Corollary: optimal $\eff^*$ supported on $\{0,1\}$. \\ \bigskip
\onslide<1->{
For $p \in [0,1]$, let $\eff_{p} = p\delta_1 + (1-p)\delta_0$. \\ \bigskip
}
\onslide<1->{
In this case $w_{\eff_p}(0) = 0$ and $w_{\eff_p}(1) = p^{\lambda/\rho}$, so optimal $p^*$ solves
\[
\max_{p \in [0,1]} \left[p - p^{\frac{\lambda + \rho}{\rho}}\right].
\]
}
\onslide<1->{
Again, objective is concave, so FOC gives the solution:
\[
p^* = \left[\frac{\rho/\lambda }{1 + \rho/\lambda} \right]^{\rho/\lambda}.
\]
}
\end{frame}


\begin{frame}
	\frametitle{}
	\begin{center}
		{\color{darkblue} {\huge Non-Smooth Costs}}
	\end{center}
\end{frame}



\begin{frame}{Non-Smooth Costs}
\textbf{Maintain:} $C$ is convex and continuous. \\ \bigskip
\textbf{Drop:} $C$ is Gateaux differentiable. \\ \bigskip

In this case, can use \textbf{directional derivative}:

\[
\begin{split}
d^{+}_{\eff}C: & \Delta X \rightarrow \mathbb{R} \cup \{\infty\},\\
\effb & \mapsto \lim_{\epsilon\searrow0}\frac{1}{\epsilon}\left[C(\eff+\epsilon(\effb - \eff)) - C(\eff) \right].
\end{split}
\]
(convexity of $C$ ensures above limit is well-defined).

\end{frame}

\begin{frame}{Non-Differentiable First-Order Approach}
\textbf{Lemma*.} For a bounded $v:X\rightarrow \mathbb{R}$, and $\eff \in 
\Eff$, 
\begin{equation*}
\eff \in \arg \max_{\effb\in \Eff}\left[ \int v(x)\effb\left( \dd x\right) -C(\effb)\right]
\end{equation*}%
if and only if 
\begin{equation*}
\eff \in \arg \max_{\effb \in \Eff}\left[ \int v(x) \effb(\dd x) - d^{+}_{\eff}C(\effb) \right]
\end{equation*}%
\onslide<2->{
Unfortunately, not nearly as useful without differentiability. 
}
\end{frame}


\begin{frame}{Implementability}

Say $\eff$ is \textbf{implementable} if $(w,\eff)$ is IC for some $w \in W$. \\ \bigskip

\textbf{Theorem 2.} The implementable distributions are dense in $\Delta X$.\\ \bigskip
\onslide<2->{
\textbf{Proof.}  \\ \medskip
If $X$ is finite, $C$ is differentiable almost everywhere. \\ \medskip 
}\onslide<3->{
For infinite $X$, show: \medskip
\begin{itemize}
\item Embed $\Delta X$ in $L^{2}(X)$. \\ \medskip
\item Show $\eff$ is implementable if $\partial_{L^2} C(\eff) \neq \varnothing$.\\ \medskip
\item Bronsted-Rockafellar: $\partial_{L^2} C(\eff) \neq \varnothing$ holds for a dense set.
\end{itemize}
}
\end{frame}

\begin{frame}{Monotonicity}
\textbf{Theorem 3.} Suppose $C$ is FOSD monotone. If $(w,\eff)$ is IC, then
\[
\bar{w}(x) = \sup_{y \leq x} w(y)
\]
is monotone, $\bar{w}$ = $w$ $\eff$-almost surely, and $(\bar{w},\eff)$ is IC. \\ \bigskip
\end{frame}

\begin{frame}{Montonicity Proof Sketch}
Suppose $X$ is finite. Every $\effb$ admits a $\effc_{\effb}$ s.t. $\effb \succeq_{FOSD} \effc_{\effb}$ and 
\[
\int \bar{w}(x) \effb(\dd x) = \int w(x) \effc_{\effb}(\dd x).
\]
\onslide<2->{
Therefore, every $\effb$ has
}
\[
\begin{split}
\onslide<2->{
\int \bar{w}(x) \only<2-5>{\effb}\only<6>{\textcolor{blue}{\eff}} (\dd x) - C(\only<2-5>{\effb}\only<6>{\textcolor{blue}{\eff}}) 
& \leq \int w(x) \effc_{\only<2-5>{\effb}\only<6>{\textcolor{blue}{\eff}}} (\dd x) - C(\effc_{\only<2-5>{\effb}\only<6>{\textcolor{blue}{\eff}}}) 
}
\onslide<3->{
\\ & \leq \int w(x) \eff (\dd x) - C(\eff) 
}
\onslide<4->{
\leq \int \bar{w}(x) \eff (\dd x) - C(\eff).
}
\end{split}
\]
\onslide<5->{So, $(\bar{w},\eff)$ is IC.}
\onslide<6->{Setting $\effb = \eff$ above gives $\bar{w}=w$ $\eff$-almost surely.\\}
\onslide<3>{(because $(w,\eff)$ is IC)}
\end{frame}

\begin{frame}
	\frametitle{}
	\begin{center}
		{\color{darkblue} {\huge Conclusion}}
	\end{center}
\end{frame}


\begin{frame}{Related Literature}
\begin{itemize}
\item \textbf{Less general flexible models}: Holmstrom and Milgrom (1987), Diamond (1998), Mirrlees and Zhou (2006), Hebert (2018), Bonham (2021), Mattsson and Weibull (2022), Bonham and Riggs-Cragun (2023). \bigskip

\item \textbf{Flexible Monitoring:} Georgiadis and Szentes (2020), Mahzoon, Shourideh, and Zetlin-Joines (2022), Wong (2023). \bigskip

\item \textbf{Robust contracting:} Carroll (2015), Antic (2022), Antic and Georgiadis (2022), Carroll and Walton (2022).
\end{itemize}
\end{frame}

\begin{frame}{Flexible Moral Hazard Problems}
We showed that in smooth \& flexible moral hazard problems:\medskip
\begin{itemize}
\item Parameters driving contract: $k_\alpha$ and $u$.\medskip
\item Cost minimization is trivial. \medskip
\item Every distribution can be implemented. \medskip
\item Wages are monotone without loss. \bigskip
\end{itemize}
\onslide<2->{
Also obtained results about principal optimality.\medskip
\begin{itemize}
\item First order approach is valid. \medskip
\item Optimality of single, binary, and discrete distributions. 
\end{itemize}
}
\end{frame}


\begin{frame}
	\frametitle{}
	\begin{center}
		{\color{darkblue} {\huge Thanks!}}
	\end{center}
\end{frame}


\end{document}
